\subsection{Conclusions and Future Directions}

This study re-evaluates the methodological prerequisites for molecular property prediction in low-data environments ($N = \NCompounds$), yielding three principal findings. First, classical machine learning ensembles operating on 2D topological fingerprints provide sufficient explicit inductive bias ($R^2 = \MorganRTwoMean \pm \MorganRTwoStd$) to surpass the predictive utility of unaugmented deep learning architectures ($R^2 = \ChemPropDMPNNFoldRTwoMean$) and computationally expensive 3D descriptors. This advantage, however, remains strictly bounded to low-data regimes where complex neural architectures suffer from parameter underdetermination. Second, rigorous nested cross-validation remains imperative to eliminate optimistic performance bias and accurately quantify out-of-sample generalization. 

Third, integrating ligand-based statistical predictions with physics-based molecular docking and ADMET profiling establishes a robust consensus filter. This orthogonal validation mitigates the false-positive rates inherent to isolated screening methodologies. By enforcing strict pharmacokinetic constraints—specifically evaluating hERG cardiotoxicity and Caco-2 intestinal permeability—the pipeline isolates \DockNValidated{} structurally and translationally viable candidates. Most notably, \DockBestCandidate{} emerges as a high-priority candidate for therapeutic repurposing, demonstrating both strong predicted affinity and a compliant safety profile. 

Future investigations must transition these computational findings into experimental validation. \textit{In vitro} enzymatic assays are required to confirm the inhibitory kinetics of \DockBestCandidate{} and the remaining consensus candidates against TgDHFR. Subsequent \textit{in vivo} efficacy models in \textit{Toxoplasma gondii}-infected cohorts will determine their systemic viability. Additionally, expanding the computational framework to evaluate multi-target selectivity profiles—specifically quantifying the binding affinity differential between TgDHFR and human DHFR—will be necessary to anticipate and minimize off-target toxicity in clinical applications.

\section*{Data and Code Availability}

All source code, Python scripts for dataset processing, nested cross-validation, model training, and docking validation protocols, as well as the benchmark datasets, are publicly available in the GitHub repository at the following \href{https://github.com/aavzpy2023/analysispaperKarimova2025}{\underline{link}}

\section*{Declaration of Generative AI's use}

During the preparation of this work, the author(s) utilized generative artificial intelligence technologies (Large Language Models) solely to assist with language polishing, structural formatting, and improving the overall readability of the manuscript. After using these tools, the author(s) thoroughly reviewed, edited, and validated all generated text. The author(s) take full and sole responsibility for the final content, scientific integrity, and accuracy of this publication. No AI technologies were used to generate data, perform statistical analyses, or draw scientific conclusions.